\section{The Relation-Induced Context}
\label{sec:contexts}
We now give the mathematical construction independently of the neural
model. This preserves the proof pattern of the companion paper and
makes explicit which properties come from order theory rather than from
training. Throughout this section, $\U$ can be any complete lattice.

\subsection{From individual summaries to a Galois ideal}
\label{sub:galois-ideal}
Let $X$ be a set and $J\subseteq X\times\U$ any relation. Assign each
object the summary
\begin{equation}\label{eq:relation-summary}
      u_{x}=\bigwedge\{u\in\U\mid(x,u)\in J\}.
\end{equation}
The empty meet is $\top_{\U}$. For visual data, $J$ is the graph of the
chosen summary map, so this meet is over a singleton; it does not replace
the spatial maximum in \cref{eq:max-summary}. Put $\X=\mathcal{P}(X)$ and
\begin{equation}\label{eq:incidence}
      (A,u)\in I\quad\Longleftrightarrow\quad
      u\leq u_{x}\text{ for every }x\in A.
\end{equation}
Here $A,B\in\X$ denote collections of objects and $u,v\in\U$ denote
attribute descriptions. Individual images are indexed by $x$.

\begin{definition}[Galois ideal and lattice-theoretic context]
      \label{def:galois-ideal}
      For complete lattices $\X$ and $\U$, a relation
      $I\subseteq\X\times\U$ is a Galois ideal if it is a lower set in the
      product order and is closed under arbitrary joins in each coordinate
      separately. Thus, including empty families,
      \begin{equation}\label{eq:ideal-joins}
            \begin{aligned}
                  (A_{t},u)\in I\ (t\in T) & \Longrightarrow
                  (\bigvee_{t\in T}A_{t},u)\in I,            \\
                  (A,u_{t})\in I\ (t\in T) & \Longrightarrow
                  (A,\bigvee_{t\in T}u_{t})\in I.
            \end{aligned}
      \end{equation}
      The triple $(\X,\U,I)$ is then a lattice-theoretic formal context
      \cite{FCA_Latt_Theor}.
\end{definition}

\begin{proposition}[Relation-induced Galois ideal]
      \label{prop:galois-ideal}
      The relation in \cref{eq:incidence} is a Galois ideal.
\end{proposition}
\begin{proof}
      Suppose $(A,u)\in I$, $B\subseteq A$, and $v\leq u$. For every $x\in B$,
      we have $v\leq u\leq u_{x}$, so $(B,v)\in I$. This proves the lower-set
      property. If $(A,u_{t})\in I$ for every $t\in T$, each $u_{x}$ with
      $x\in A$ is an upper bound of the family, hence
      $\bigvee_{t\in T}u_{t}\leq u_{x}$. If $(A_{t},u)\in I$ for every
      $t\in T$, every $x\in\bigcup_{t\in T}A_{t}$ belongs to some $A_{t}$ and
      satisfies $u\leq u_{x}$. The first assertion also holds for an empty
      family because its join is $\bot_{\U}$; the second is then vacuous.
\end{proof}

\subsection{Derivations, closure, and concepts}
\label{sub:derivations}
The summary and retrieval maps are
\begin{equation}\label{eq:derivations}
      F_{I}(A)=\bigwedge_{x\in A}u_{x},\qquad
      G_{I}(u)=\{x\in X\mid u\leq u_{x}\}.
\end{equation}
Thus $F_{I}(\varnothing)=\top_{\U}$ and $G_{I}(\bot_{\U})=X$.

\begin{proposition}[Summary--retrieval adjunction]
      \label{prop:adjunction}
      The maps $F_{I}:\X\to\U$ and $G_{I}:\U\to\X$ are antitone and satisfy
      \begin{equation}\label{eq:adjunction}
            A\subseteq G_{I}(u)\quad\Longleftrightarrow\quad u\leq F_{I}(A).
      \end{equation}
      Consequently, $c_{I}=G_{I}F_{I}$ and $d_{I}=F_{I}G_{I}$ are closure
      operators on $\X$ and $\U$, respectively.
\end{proposition}
\begin{proof}
      If $B\subseteq A$, taking a meet over the larger set gives
      $F_{I}(A)\leq F_{I}(B)$. If $u\leq v$, every summary dominating $v$
      also dominates $u$, so $G_{I}(v)\subseteq G_{I}(u)$. Finally,
      $A\subseteq G_{I}(u)$ means exactly that $u\leq u_{x}$ for all
      $x\in A$, equivalently $u\leq\bigwedge_{x\in A}u_{x}$.

      The equivalence gives $A\subseteq G_{I}F_{I}(A)$ and
      $u\leq F_{I}G_{I}(u)$. The composites are monotone since each factor is
      antitone. Applying antitonicity to these inequalities gives
      $F_{I}G_{I}F_{I}=F_{I}$ and $G_{I}F_{I}G_{I}=G_{I}$, hence idempotence
      of both composites. These are precisely extensivity, monotonicity, and
      idempotence, the defining properties of closure operators.
\end{proof}

\begin{proposition}[Identification of context derivations]
      \label{prop:identification}
      The lattice-theoretic derivations associated with $I$ are exactly
      $A^{I}=F_{I}(A)$ and $u^{I}=G_{I}(u)$.
\end{proposition}
\begin{proof}
      By definition, $u^{I}$ is the join of all $A$ incident with $u$, here
      $\bigcup\{A\in\X\mid(A,u)\in I\}$. Every such $A$ is contained in
      $G_{I}(u)$; conversely, each $x\in G_{I}(u)$ supplies the incident set
      $\{x\}$. Thus this union is $G_{I}(u)$. Similarly, $A^{I}$ is the join
      of all $u$ incident with $A$. Every such $u$ lies below $F_{I}(A)$,
      and $F_{I}(A)$ is itself incident with $A$, proving equality.
\end{proof}

A lattice-theoretic formal concept is a pair $(A,u)$ satisfying
$F_{I}(A)=u$ and $G_{I}(u)=A$. The set $A$ is its \emph{extent}, and $u$
is its \emph{intent}. Order the concepts by
\begin{equation}\label{eq:concept-order}
      (A,u)\leq(B,v)\quad\Longleftrightarrow\quad
      A\subseteq B\quad\Longleftrightarrow\quad v\leq u.
\end{equation}
A larger extent has a weaker shared description. Closed concepts record
all consequences of a query on the chosen dataset, not on unseen images.

\begin{theorem}[Concept lattice and generated concepts]
      \label{thm:concept-lattice}
      The concepts form a complete lattice $\B_{I}$. For any indexed family
      $\{(A_{t},u_{t})\}_{t\in T}$,
      \begin{equation}\label{eq:concept-operations}
            \begin{aligned}
                  \bigwedge_{t\in T}(A_{t},u_{t})
                   & =\left(\bigcap_{t\in T}A_{t},
                  d_{I}(\bigvee_{t\in T}u_{t})\right),    \\
                  \bigvee_{t\in T}(A_{t},u_{t})
                   & =\left(c_{I}(\bigcup_{t\in T}A_{t}),
                  \bigwedge_{t\in T}u_{t}\right).
            \end{aligned}
      \end{equation}
      The generated concepts are
      \begin{equation}\label{eq:generated-concepts}
            \gamma(A)=(c_{I}(A),F_{I}(A)),\qquad
            \mu(u)=(G_{I}(u),d_{I}(u)).
      \end{equation}
\end{theorem}
\begin{proof}
      Closed sets of a closure operator are closed under arbitrary
      intersections: if each $A_{t}$ is closed, monotonicity gives
      $c_{I}(\bigcap_{t}A_{t})\subseteq A_{t}$ for every $t$, and extensivity
      gives the reverse inclusion to their intersection. Their least closed
      upper bound is the closure of their union. Also, directly from the
      adjunction,
      \begin{equation}\label{eq:derivation-family}
            G_{I}(\bigvee_{t\in T}u_{t})=\bigcap_{t\in T}G_{I}(u_{t}),\qquad
            F_{I}(\bigcup_{t\in T}A_{t})=\bigwedge_{t\in T}F_{I}(A_{t}).
      \end{equation}
      Together with $F_{I}G_{I}F_{I}=F_{I}$, these identities give the stated
      intents and establish both formulas, including empty families. The same
      identities prove that $\gamma(A)$ and $\mu(u)$ are concepts.
\end{proof}

The maps in \cref{eq:generated-concepts} need not be embeddings: distinct
queries or source sets can generate the same concept. In particular,
$G_{I}(d_{I}(u))=G_{I}(u)$. A join of raw queries conjoins requirements,
whereas a meet relaxes them:
\begin{equation}\label{eq:query-operations}
      G_{I}(u\vee v)=G_{I}(u)\cap G_{I}(v),\qquad
      G_{I}(u\wedge v)\supseteq G_{I}(u)\cup G_{I}(v).
\end{equation}
The inclusion can be strict. These operations on queries should not be
confused with the extent-ordered concept operations in
\cref{eq:concept-operations}.

\subsection{A worked example and classical representation}
\label{sub:example-scaling}
Suppose three images have summaries
\begin{equation}\label{eq:worked-summaries}
      u_{x_{1}}=(4,1),\qquad u_{x_{2}}=(2,3),\qquad u_{x_{3}}=(1,0).
\end{equation}
Then $\U=[0,4]\times[0,3]$. The query $(2,1)$ retrieves $x_{1},x_{2}$
and is already closed. The query $(3,1)$ retrieves only $x_{1}$ and
closes to $(4,1)$. Both queries have the same positive support, so a
presence-only representation cannot distinguish them. In the other
direction, the singleton $\{x_{3}\}$ closes to all three images because
every summary dominates $(1,0)$. Finally,
\begin{equation}\label{eq:worked-meet-join}
      \begin{aligned}
            G_{I}((3,0)\vee(0,2))   & =\varnothing, \\
            G_{I}((3,0)\wedge(0,2)) & =X
            \supsetneq\{x_{1},x_{2}\}.
      \end{aligned}
\end{equation}
Here a stronger conjunction is empty, whereas the shared lower bound
is a completely unselective query. Neither outcome is mathematically
exceptional.

For a functional summary map, this construction is a pattern structure:
each object has a description in a meet-semilattice, and common
descriptions are obtained by meets~\cite{ganter2001pattern}. It is also
representable by classical FCA without discarding magnitudes.

\begin{proposition}[Observed-threshold representation]
      \label{prop:ordinal-scaling}
      Let $X$ be finite and nonempty, with nonnegative real summaries and $\U$
      as in \cref{eq:vector-lattice}. Define classical attributes
      \begin{equation}\label{eq:ordinal-attributes}
            E=\{(j,a)\mid 1\leq j\leq d,\ a>0,
            \ a=u_{x,j}\text{ for some }x\in X\}.
      \end{equation}
      Declare $x$ incident with $(j,a)$ exactly when $u_{x,j}\geq a$. This
      classical context and the vector context have identical closed extents;
      their concept lattices are therefore isomorphic by preservation of
      extents.
\end{proposition}
\begin{proof}
      The classical intent of $A\subseteq X$ consists precisely of the pairs
      $(j,a)$ with $a\leq F_{I}(A)_{j}$. Each coordinate of $F_{I}(A)$ is an
      observed value or zero, including $A=\varnothing$, whose meet is $M$.
      Thus satisfying every attribute in that classical intent is equivalent
      to dominating $F_{I}(A)$. Consequently $A''=G_{I}F_{I}(A)$. Since a
      concept is determined by its extent and ordered by extent inclusion,
      the equality of extent closures proves the isomorphism.
\end{proof}

An arbitrary $u\in\U$ can be rounded coordinatewise upward to the next
observed value without changing its extent. Presence-only FCA retains
just one bit per coordinate, whereas ordinal scaling retains the relevant
thresholds. Any advantage over a presence baseline is therefore an
advantage of retaining amplitudes, not an expressive advantage over FCA.
The vector representation gives these amplitudes and their operations
directly, without materializing a larger binary incidence matrix.
